\documentclass{article}
\usepackage{amsmath}
\usepackage{amssymb}
\begin{document}

\title{Glossary of Set Notation Symbols}
\author{}
\date{}
\maketitle

\section{Set Operations Glossary}

\begin{itemize}
    \item \(\emptyset\): The **empty set**, the set with no elements.
    \[
    \emptyset = \{\}
    \]
    \item \(\in\): **Element of**. Denotes that an element is part of a set.
    \[
    x \in A \quad \text{means} \quad x \text{ is an element of set } A
    \]
    \item \(\notin\): **Not an element of**. Denotes that an element is not part of a set.
    \[
    x \notin A \quad \text{means} \quad x \text{ is not an element of set } A
    \]
    \item \(\subseteq\): **Subset**. Set \( A \) is a subset of set \( B \) if all elements of \( A \) are also elements of \( B \).
    \[
    A \subseteq B \quad \text{means} \quad \text{every element of } A \text{ is also in } B
    \]
    \item \(\subset\): **Proper subset**. Set \( A \) is a proper subset of \( B \) if all elements of \( A \) are in \( B \), but \( A \neq B \).
    \[
    A \subset B \quad \text{means} \quad A \text{ is a proper subset of } B
    \]
    \item \(\cup\): **Union**. The union of sets \( A \) and \( B \) is the set of elements that are in \( A \), \( B \), or both.
    \[
    A \cup B = \{x \mid x \in A \text{ or } x \in B\}
    \]
    \item \(\cap\): **Intersection**. The intersection of sets \( A \) and \( B \) is the set of elements that are in both \( A \) and \( B \).
    \[
    A \cap B = \{x \mid x \in A \text{ and } x \in B\}
    \]
    \item \(\setminus\): **Set difference**. The difference between sets \( A \) and \( B \) is the set of elements that are in \( A \) but not in \( B \).
    \[
    A \setminus B = \{x \mid x \in A \text{ and } x \notin B\}
    \]
    \item \(\triangle\): **Symmetric difference**. The symmetric difference of sets \( A \) and \( B \) is the set of elements that are in either \( A \) or \( B \), but not in both.
    \[
    A \triangle B = (A \setminus B) \cup (B \setminus A)
    \]
    \item \(\mathbb{U}\): **Universal set**. The set containing all possible elements under consideration in a particular discussion.
    \[
    A \subseteq \mathbb{U} \quad \text{where } \mathbb{U} \text{ is the universal set.}
    \]
    \item \(A^c\): **Complement**. The complement of set \( A \) is the set of all elements in the universal set \( \mathbb{U} \) that are not in \( A \).
    \[
    A^c = \{x \mid x \in \mathbb{U} \text{ and } x \notin A\}
    \]
    \item \(|A|\): **Cardinality**. The cardinality of a set is the number of elements in the set \( A \).
    \[
    |A| = \text{number of elements in } A
    \]
\end{itemize}

\end{document}
